\section{Background}

Consider the space of vector-valued functions $f:\{-1,1\}^n \rightarrow X$, defined on the discrete hypercube, for which $(X,\|\cdot\|)$ is a normed space.  Define the $L^p(X)$ norm on such functions as 
\[\|f\|_{p}  = \left(\mathbb{E}\|f\|^p \right)^{1/p} = \Big(\frac{1}{2^n}\sum\limits_{\varepsilon \in \{-1,1\}^n} \|f(\varepsilon)\|^p \Big)^{1/p}.\]
It is often helpful to think about $f$ in terms of its  \textit{Fourier-Walsh} expansion
\[f(\varepsilon) = \sum_{S \subset \{1,...,n\}} a_S\varepsilon^S,\]
where 
\[\varepsilon^S: (\varepsilon_1,...,\varepsilon_n) \longmapsto \prod\limits_{i \in S} \varepsilon_{i}, \quad \varepsilon^{\emptyset} \equiv 1,\]
are called \textit{Walsh functions}.  A  straightforward computation yields 
\[a_S = \mathbb{E}\varepsilon^{S}f(\varepsilon) \quad \text{ for all } \quad S\subseteq \{1,...,n\}.\]
Much of what we will prove in this paper relies on the rich theory of the \textit{heat semigroup} on the discrete hypercube.  With the discrete derivative operators $D_j$ defined as in the introduction, define the \textit{discrete Laplacian} $\Delta$ on functions $f:\{-1,1\}^n \to X$ as follows
\[\Delta f(\varepsilon)= - \sum_{j=1}^{n} D_{j} f(\varepsilon).\]
We then define the \textit{heat semigroup} $P_t = e^{\Delta t}$ as
\[P_{t}(f)(\varepsilon) = \sum_{S \subset \{1,...,n\}} a_Se^{-|S|t}\varepsilon^S \quad \text{for all}\quad t \geq 0.\]
An important property satisfied by the heat semigroup is \textit{hypercontractivity}: for all $p,q$ satisfying $1 < p \leq q < \infty$ and $e^{-2t} \leq \frac{p-1}{q-1}$, we have   
\begin{align}\label{hyp1}
\|P_t f\|_{q} \leq \|f\|_{p}.
\end{align}


One of the key ingredients in obtaining Theorem~\ref{mth1}  is going to be the following pointwise identity obtained in \cite{IVHV}.  For a normed space $(X,||\cdot||)$ and a function $f:\{-1,1\}^n \to X$, we have 
\begin{equation}
    -\frac{d}{dt}P_t f(\varepsilon) = \frac{1}{\sqrt{e^{2t} - 1}}\mathbb{E}_\xi \sum_{j=1}^n \delta_j(t) D_j f(\varepsilon\xi(t)), \label{heat}
\end{equation}
where $\varepsilon\xi(t) = (\varepsilon_1\xi_1(t),...,\varepsilon_n\xi_n(t))$, and the $\xi_i(t)$ are i.i.d. random variables with
\[\mathbb{P}\{\xi_i(t) = \pm 1\} = \frac{1\pm e^{-t}}{2},\]
and   $\delta_i = \frac{\xi_i(t) - \mathbb{E} \xi_{i}(t)}{\sqrt{\mathrm{Var}(\xi_{i}(t))}}$.  
